
\documentclass[12pt]{article}
\usepackage[utf8]{inputenc}
\usepackage[russian]{babel}
\usepackage{amsmath}
\usepackage{amssymb}
\usepackage{geometry}
\geometry{a4paper, margin=2cm}

\begin{document}

\section*{Задача 1.2}

\textbf{Условие:} Написать программу прибавления $1$ к двоичному
числу, записанному правильно. Провести анализ в среднем и в худшем
случаях.

\section*{Программа}

\begin{verbatim}
alphabet = [#, 0, 1]
tapes = [ right ]

START:
    # -> #R CheckEmpty
CheckEmpty:
    # -> #H STOP
    1 -> 1R RMove
    0 -> 0R RMove
RMove:
    # -> #L Inc
    1 -> 1R RMove
    0 -> 0R RMove
Inc:
    0 -> 1H STOP
    1 -> 0L Inc
    # -> 1H ShiftStart
ShiftStart:
    1 -> #R Shift_1
Shift_1:
    1 -> 1R Shift_1
    0 -> 1R Shift_0
    # -> 1H STOP
Shift_0:
    1 -> 0R Shift_1
    0 -> 0R Shift_0
    # -> 0H STOP
\end{verbatim}

\subsection*{Сложность}

Пусть $n = |w|$ без $\#$. Пусть в $w$ содержится хотя бы один ноль. Пусть $w=u01^i$. $u=\varnothing$ если $w=0$

До входа в $Inc$:
\begin{enumerate}
    \item $START$ - один шаг
    \item $CheckEmpty$ - один шаг
    \item $RMove$: проход по оставшимся цифрам вправо и переход на правом $\#$ - $n$ шагов
\end{enumerate}
Итого до $Inc$ - $n+2$ переходов

$Inc$: каждую из $i$ единиц на меняем ноль, затем первый встретившийся ноль меняется на единицу и останавливается. Это занимает $i+1$ шагов. Следовательно, для $w=u01^i$: 
\[
    T(n,i)=n+i+3
\]

Отдельно есть случай $\#0\#$. Тогда $T(\#0\#)=4$.

Теперь рассмотрим $w=1^n$. Первые состояния до $Inc$ занимают также $n+2$ шага. В состоянии $Inc$ машина проходит влево по всем единицам и затем обрабатывает левый ограничитель - $n+1$ шаг. Потом:
\begin{enumerate}
    \item $ShiftStart$ - 1 шаг
    \item $Shift\_1$ - 1 шаг
    \item $Shift\_0$: сдвиг оставшихся $n-1$ цифр и плюс обработка правого ограничителя - $n$ шагов
\end{enumerate}
В итоге:
\[
T(1^n)=(n+2)+(n+1)+1+1+n = 3n+5
\]

\subsection*{Сложность в худшем}

Если вход содержит ноль, то $T \le 2n+1$. Вход из одних единиц - $3n+5$ шагов. Значит худший случай - $w=1^n$.

Значит сложность в худшем случае:
\[
T_{\max}(n)=3n+5
\]

\subsection*{Сложность в среднем}

Случай $\#0\#$ отдельный. Значит рассмотрим для $n=1$:
\[
T_{\mathrm{avg}}(1)=\frac{4+8}{2}=6
\]

Всего слов длины $n$ - $2^{n-1}$. Слов $1u01^i$ - $2^{n-i-2}$, значит вероятность: $\frac{2^{n-i-2}}{2^{n-1}}=\frac{1}{2^{i+1}}$

Ну и плюс случай $1^n$ с вероятностью $\frac{1}{2^{n-1}}$. Тогда среднее время работы:
\[
T_{avg}(n)
=\sum_{i=0}^{n-2}\frac{n+i+3}{2^{i+1}}
 +\frac{3n+5}{2^{n-1}}
\]
По WolframAlpha:
\[
\sum_{i=0}^{n-2}\frac{1}{2^{i+1}}=1-\frac{1}{2^{n-1}},
\sum_{i=0}^{n-2}\frac{i}{2^{i+1}}=1-\frac{n}{2^{n-1}}
\]
Тогда
\[
T_{avg}(n)
=(n+3)\left(1-\frac{1}{2^{n-1}}\right)
  +1-\frac{n}{2^{n-1}}+\frac{3n+5}{2^{n-1}}=n+4+\frac{n+2}{2^{n-1}}
\]
Значит, сложность в среднем:
\[
T_{avg}(n)=
\begin{cases}
6, & n=1\\
 n+4+\dfrac{n+2}{2^{n-1}}, & n\geq2
\end{cases}
\]

\end{document}